% part: intuitionistic-logic
% chapter: soundness-completeness
% section: canonical-model

\documentclass[../../../include/open-logic-section]{subfiles}

\begin{document}

\olfileid[hi]{int}{sc}{mod}

\olsection{विहित प्रतिरूप}

हमारे प्रतिरूप के जगत प्राकृतिक संख्याओं के परिमित अनुक्रम~$\sigma$ होंगे;
अर्थात $\sigma \in \Nat^*$। ध्यान दें कि $\Nat^*$ आगमनात्मक रूप से इस
प्रकार परिभाषित है:
\begin{enumerate}
\item $\emptyseq \in \Nat^*$.
\item यदि $\sigma \in \Nat^*$ और $n \in \Nat$, तो
  $\sigma.n \in \Nat^*$ (जहाँ $\sigma.n$, $\sigma \concat \tuple{n}$ है,
  और $\sigma \concat \sigma'$, $\sigma$ तथा $\sigma'$ का सन्निवेशन है)।
\item इसके अतिरिक्त कुछ भी $\Nat^*$ में नहीं है।
\end{enumerate}
अतः हम आगमनात्मक परिभाषाएँ देने के लिए $\Nat^*$ का उपयोग कर सकते हैं।

मान लें कि $\tuple{!B_1, !C_1}$, $\tuple{!B_2, !C_2}$, \dots,
!!{formula}ों के सभी युग्मों का परिगणन है। !!{formula}ों का समुच्चय~$\Delta$
दिए जाने पर $\Delta(\sigma)$ को आगमन द्वारा इस प्रकार परिभाषित करें:
\begin{enumerate}
\item $\Delta(\emptyseq) = \Delta$
\item $\Delta(\sigma.n) = {}$
  \[
  \begin{cases}
    (\Delta(\sigma) \cup \{!B_n\})^* &
    \text{यदि $\Delta(\sigma) \cup \{!B_n\} \Proves/ !C_n$} \\
    \Delta(\sigma) & \text{अन्यथा}
  \end{cases}
  \]
\end{enumerate}
यहाँ $(\Delta(\sigma) \cup \{!B_n\})^*$ से हमारा आशय !!{formula}ों के उस
अभाज्य समुच्चय से है जिसका अस्तित्व समुच्चय
$\Delta(\sigma) \cup \{!B_n\}$ और !!{formula}~$!C_n$ पर
\olref[lin]{lem:lindenbaum} लगाने से मिलता है। ध्यान दें कि इस परिभाषा से,
यदि $\Delta(\sigma) \cup \{!B_n\} \Proves/ !C_n$, तो
$\Delta(\sigma.n) \Proves !B_n$ और
$\Delta(\sigma.n) \Proves/ !C_n$। यह भी ध्यान दें कि किसी भी~$n$ के लिए
$\Delta(\sigma) \subseteq \Delta(\sigma.n)$। यदि $\Delta$ अभाज्य है, तो
प्रत्येक~$\sigma$ के लिए $\Delta(\sigma)$ अभाज्य है।

\begin{defn}\ollabel{defn:canonical-model}
  मान लें कि $\Delta$ अभाज्य है। तब $\Delta$ का \emph{विहित प्रतिरूप}
  $\mModel{M(\Delta)}$ इस प्रकार परिभाषित है:
  \begin{enumerate}
  \item $W = \Nat^*$, प्राकृतिक संख्याओं के परिमित अनुक्रमों का समुच्चय।
  \item $R$ वह आंशिक क्रम है जिसके अनुसार $R\sigma\sigma'$ तभी जब $\sigma$,
    ~$\sigma'$ का आरंभिक खंड हो (अर्थात किसी अनुक्रम~$\sigma''$ के लिए
    $\sigma' = \sigma \concat \sigma''$)।
  \item $V(p) = \Setabs{\sigma}{p \in \Delta(\sigma)}$.
  \end{enumerate}
\end{defn}

यह सत्यापित करना सरल है कि $R$ वास्तव में आंशिक क्रम है।~$V$ पर एकदिष्टता
की शर्त भी संतुष्ट होती है। चूँकि
$\Delta(\sigma) \subseteq \Delta(\sigma.n)$, इसलिए~$\sigma$ पर आगमन द्वारा
जब भी $R\sigma\sigma'$, तब
$\Delta(\sigma) \subseteq \Delta(\sigma')$ मिलता है।

\end{document}
